\section{Creando Nuestras Propias Herramientas - Funciones}\n\\label{sec:funciones}\n\nEn los capítulos anteriores, has aprendido a declarar variables, operar con ellas, tomar decisiones y repetir tareas. Pero a medida que tus programas crecen en complejidad, mantener todo el código en la función \\texttt{main()} se vuelve cada vez más difícil y propenso a errores. Imagina si tuvieras que construir una casa entera sin dividirla en habitaciones: sería caótico, difícil de mantener y casi imposible de mejorar.\n\nLas \\textbf{funciones} son como habitaciones en tu programa: dividen tu código en piezas manejables, reutilizables y fáciles de entender. Son una de las herramientas más poderosas del programador moderno.\n\n\\subsection{El Principio DRY: No Repitas Código}\n\\label{subsec:principio_dry}\n\nAntes de aprender a escribir funciones, debes comprender un principio fundamental de la programación: \\textbf{DRY (Don't Repeat Yourself)}. Este principio dice que cada parte del conocimiento en un sistema debe tener una representación única, no ambigua y definitiva dentro del sistema.\n\n¿Por qué es importante? Miremos un ejemplo sin funciones:\n\n\\begin{lstlisting}[language=C++]\n#include <iostream>\n#include <cmath>\nusing namespace std;\n\nint main() {\n    // Calcular area de un circulo (repeticion)\n    double radio1 = 5.0;\n    double area1 = 3.14159 * radio1 * radio1;\n    cout << \"Area del circulo 1: \" << area1 << endl;\n    \n    // Calcular area de otro circulo (repeticion)\n    double radio2 = 7.5;\n    double area2 = 3.14159 * radio2 * radio2;\n    cout << \"Area del circulo 2: \" << area2 << endl;\n    \n    // Calcular area de otro circulo (repeticion)\n    double radio3 = 3.2;\n    double area3 = 3.14159 * radio3 * radio3;\n    cout << \"Area del circulo 3: \" << area3 << endl;\n    \n    return 0;\n}\n\\end{lstlisting}\n\n¿Ves la repetición? Si necesitamos cambiar la fórmula (por ejemplo, usar un valor más preciso de Pi), tendríamos que cambiarla en múltiples lugares. Esto es propenso a errores y difícil de mantener.\n\nAhora veamos cómo se ve con una función:\n\n\\begin{lstlisting}[language=C++]\n#include <iostream>\n#include <cmath>\nusing namespace std;\n\n// Definicion de la funcion para calcular area de circulo\ndouble calcularAreaCirculo(double radio) {\n    return 3.14159 * radio * radio;\n}\n\nint main() {\n    // Usar la funcion (no hay repeticion!)\n    cout << \"Area del circulo 1: \" << calcularAreaCirculo(5.0) << endl;\n    cout << \"Area del circulo 2: \" << calcularAreaCirculo(7.5) << endl;\n    cout << \"Area del circulo 3: \" << calcularAreaCirculo(3.2) << endl;\n    \n    return 0;\n}\n\\end{lstlisting}\n\nCon la función, si necesitamos cambiar la fórmula, solo lo hacemos en un lugar. ¡Ese es el poder del principio DRY!\n\n\\subsection{Anatomía de una Función: Componentes Esenciales}\n\\label{subsec:anatomia_funcion}\n\nToda función en C++ tiene una estructura específica que debes conocer. La anatomía de una función es:\n\n\\begin{lstlisting}[language=C++]\ntipoDeRetorno nombreDeLaFuncion(tipo parametro1, tipo parametro2, ...) {\n    // Cuerpo de la funcion\n    // Instrucciones que realizan la tarea\n    return valor; // Si la funcion devuelve algo\n}\n\\end{lstlisting}\n\nDesglosemos cada componente:\n\n\\begin{itemize}\n    \\item \\textbf{Tipo de Retorno:} Especifica qué tipo de valor devuelve la función. Puede ser \\texttt{int}, \\texttt{double}, \\texttt{bool}, \\texttt{void} (nada), etc.\n    \\item \\textbf{Nombre de la Función:} El identificador que usarás para llamar a la función. Debe ser descriptivo y seguir las mismas reglas que los nombres de variables.\n    \\item \\textbf{Parámetros:} Valores que la función necesita para realizar su tarea. Son como ``ingredientes'' para la función. Pueden ser cero o más.\n    \\item \\textbf{Cuerpo de la Función:} Las instrucciones que realizan la tarea específica.\n    \\item \\textbf{Sentencia \\texttt{return}:} Devuelve un valor al código que llamó a la función y termina la ejecución de la función.\n\\end{itemize}\n\nVeamos un ejemplo completo que ilustra todos estos componentes:\n\n\\begin{lstlisting}[language=C++]\n#include <iostream>\nusing namespace std;\n\n// Funcion que calcula el area de un rectangulo\n// Tipo de retorno: double\n// Nombre: calcularAreaRectangulo  \n// Parametros: double largo, double ancho\n// Cuerpo: calcula y retorna el area\n// Sentencia return: devuelve el valor calculado\ndouble calcularAreaRectangulo(double largo, double ancho) {\n    double area = largo * ancho;\n    return area;\n}\n\nint main() {\n    // Llamada a la funcion\n    double area = calcularAreaRectangulo(10.0, 5.0);\n    cout << \"Area del rectangulo: \" << area << endl;\n    \n    return 0;\n}\n\\end{lstlisting}\n\n\\textbf{Un aspecto importante:} Las funciones deben ser definidas \\textbf{antes} de ser usadas en C++. Si defines tu función después de \\texttt{main()} y la usas en \\texttt{main()}, el compilador no sabrá que existe. Para solucionar esto, puedes declarar la función antes de usarla:\n\n\\begin{lstlisting}[language=C++]\n#include <iostream>\nusing namespace std;\n\n// Declaracion anticipada de la funcion (tambien llamada prototipo)\ndouble calcularAreaRectangulo(double largo, double ancho);\n\nint main() {\n    // Ahora podemos usar la funcion antes de que este definida\n    double area = calcularAreaRectangulo(10.0, 5.0);\n    cout << \"Area del rectangulo: \" << area << endl;\n    \n    return 0;\n}\n\n// Definicion de la funcion despues de main\ndouble calcularAreaRectangulo(double largo, double ancho) {\n    return largo * ancho;\n}\n\\end{lstlisting}\n\n\\subsection{Paso de Parámetros por Valor: Seguridad en el Traspaso}\n\\label{subsec:paso_parametros_valor}\n\nCuando pasas variables como argumentos a una función en C++, por defecto se usa el \\textbf{paso por valor}. Esto significa que la función recibe una \\textbf{copia} del valor, no la variable original.\n\nVeamos un ejemplo que ilustra esta característica:\n\n\\begin{lstlisting}[language=C++]\n#include <iostream>\nusing namespace std;\n\n// Funcion que recibe una copia del valor\nvoid duplicar(int numero) {\n    numero = numero * 2;  // Esto solo cambia la copia, no la variable original\n    cout << \"Valor dentro de la funcion: \" << numero << endl;\n}\n\nint main() {\n    int valor = 10;\n    cout << \"Valor antes de la funcion: \" << valor << endl;\n    \n    duplicar(valor);  // Pasamos una copia de 'valor'\n    \n    cout << \"Valor despues de la funcion: \" << valor << endl;  // Aun es 10!\n    \n    return 0;\n}\n\\end{lstlisting}\n\nEn este ejemplo, aunque la función \\texttt{duplicar} multiplica su parámetro por 2, el valor original de \\texttt{valor} en \\texttt{main()} no cambia. Eso es porque la función solo recibió una \\textbf{copia} del valor.\n\nEsta característica es extremadamente valiosa porque:\n\n\\begin{itemize}\n    \\item \\textbf{Seguridad:} Las funciones no pueden accidentalmente modificar tus variables originales.\n    \\item \\textbf{Claridad:} Puedes estar seguro de que pasar una variable a una función no cambiará su valor en el lugar donde la llamas.\n    \\item \\textbf{Facilidad de depuración:} Es más fácil rastrear de dónde vienen los cambios a tus variables.\n\\end{itemize}\n\n\\textbf{¿Cuándo usar paso por valor?} Es ideal cuando solo necesitas leer el valor de la variable, no modificarla. Es especialmente seguro para tipos de datos pequeños como \\texttt{int}, \\texttt{double}, \\texttt{char}, etc.\n\n\\subsection{La Sentencia \\texttt{return}: Control del Flujo de Ejecución}\n\\label{subsec:sentencia_return}\n\nLa sentencia \\texttt{return} tiene dos funciones importantes en C++:\n\n\\begin{enumerate}\n    \\item \\textbf{Devolver un valor:} Si la función tiene un tipo de retorno distinto de \\texttt{void}, \\texttt{return} envía un valor de vuelta al código que llamó a la función.\n    \\item \\textbf{Terminar la función:} \\texttt{return} inmediatamente termina la ejecución de la función, sin importar cuánto código quede por ejecutar.\n\\end{enumerate}\n\nVeamos ejemplos de ambos usos:\n\n\\begin{lstlisting}[language=C++]\n#include <iostream>\nusing namespace std;\n\n// Funcion que devuelve un valor\nint sumar(int a, int b) {\n    int resultado = a + b;\n    return resultado;  // Devuelve el valor y termina la funcion\n    // Esta linea no se ejecutara nunca:\n    cout << \"Este texto no se imprimira\" << endl;\n}\n\n// Funcion que no devuelve nada (void)\nvoid mostrarMensaje(string mensaje) {\n    cout << \"Mensaje: \" << mensaje << endl;\n    return;  // Termina la funcion, opcional en funciones void\n    // Esta linea tampoco se ejecutara:\n    cout << \"Esto tampoco se imprimira\" << endl;\n}\n\n// Uso de return para salir anticipadamente de una funcion\nbool esNumeroValido(int numero) {\n    if (numero < 0) {\n        cout << \"Error: Numero negativo no valido\" << endl;\n        return false;  // Sale inmediatamente de la funcion\n    }\n    if (numero > 100) {\n        cout << \"Error: Numero demasiado grande\" << endl;\n        return false;  // Sale inmediatamente de la funcion\n    }\n    // Si llega aqui, el numero es valido\n    return true;\n}\n\nint main() {\n    int resultado = sumar(5, 3);\n    cout << \"Resultado de la suma: \" << resultado << endl;\n    \n    mostrarMensaje(\"Hola, mundo!\");\n    \n    bool valido1 = esNumeroValido(-5);   // Devuelve false inmediatamente\n    bool valido2 = esNumeroValido(150);  // Devuelve false inmediatamente\n    bool valido3 = esNumeroValido(50);   // Devuelve true\n    \n    return 0;\n}\n\\end{lstlisting}\n\nEl uso de \\texttt{return} para salir anticipadamente de una función es una técnica muy útil para manejar errores o condiciones especiales, como se muestra en la función \\texttt{esNumeroValido}.\n\n\\subsection{Ámbito (Scope) de las Variables: Donde Viven las Variables}\n\\label{subsec:ambito_variables}\n\nEl \\textbf{ámbito} (scope) de una variable determina dónde en el código puedes acceder a esa variable. En C++, las variables tienen \\textbf{ámbito local} o \\textbf{ámbito global}.\n\n\\subsubsection{Ámbito Local}\n\nLas variables declaradas dentro de una función (incluyendo los parámetros) tienen \\textbf{ámbito local} a esa función. Solo pueden ser usadas dentro de esa función:\n\n\\begin{lstlisting}[language=C++]\n#include <iostream>\nusing namespace std;\n\nvoid funcion1() {\n    int variableLocal = 10;  // Variable local a funcion1\n    cout << \"En funcion1, variableLocal = \" << variableLocal << endl;\n    // variableLocal solo existe aqui\n}\n\nvoid funcion2() {\n    // ESTO NO FUNCIONA - variableLocal no existe aqui\n    // cout << variableLocal << endl;  // Error de compilacion\n    cout << \"En funcion2\" << endl;\n}\n\nint main() {\n    // ESTO TAMPOCO FUNCIONA - variableLocal no existe aqui\n    // cout << variableLocal << endl;  // Error de compilacion\n    \n    funcion1();\n    funcion2();\n    return 0;\n}\n\\end{lstlisting}\n\n\\subsubsection{Parámetros y Variables Locales}\n\nCada función recibe sus propios parámetros como variables locales:\n\n\\begin{lstlisting}[language=C++]\n#include <iostream>\nusing namespace std;\n\nvoid funcionConParametros(int parametro) {\n    // parametro es una variable local con el valor que se paso como argumento\n    int variableLocal = parametro * 2;  // Variable local tambien\n    \n    cout << \"parametro: \" << parametro << endl;\n    cout << \"variableLocal: \" << variableLocal << endl;\n}\n\nint main() {\n    int argumento = 5;\n    funcionConParametros(argumento);  // Paso el valor de 'argumento'\n    \n    // 'parametro' y 'variableLocal' no existen aqui\n    return 0;\n}\n\\end{lstlisting}\n\n\\subsubsection{Variables con el Mismo Nombre en Diferentes Funciones}\n\nDiferentes funciones pueden tener variables con el mismo nombre sin que se interfieran entre sí, porque cada una tiene su propio ámbito:\n\n\\begin{lstlisting}[language=C++]\n#include <iostream>\nusing namespace std;\n\nvoid funcionA() {\n    int x = 10;\n    cout << \"En funcionA, x = \" << x << endl;  // Imprime 10\n}\n\nvoid funcionB() {\n    int x = 20;  // Esta x es diferente de la de funcionA\n    cout << \"En funcionB, x = \" << x << endl;  // Imprime 20\n}\n\nint main() {\n    int x = 30;  // Esta x es diferente de las otras\n    cout << \"En main, x = \" << x << endl;  // Imprime 30\n    \n    funcionA();\n    funcionB();\n    \n    return 0;\n}\n\\end{lstlisting}\n\n\\textbf{Regla importante:} Una variable local ``oculta'' a variables con el mismo nombre en ámbitos externos. Si declaras una variable en \\texttt{main()} y otra con el mismo nombre en una función, son completamente diferentes.\n\n\\subsection{Funciones como Bloques de Construcción: Organización del Código}\n\\label{subsec:funciones_como_bloques}\n\nLas funciones no solo evitan la repetición de código, sino que también ayudan a organizar tu programa en bloques lógicos. Esto mejora la \\textbf{legibilidad}, \\textbf{mantenibilidad} y \\textbf{reutilización} del código.\n\nVeamos un ejemplo que combina todos los conceptos anteriores:\n\n\\begin{lstlisting}[language=C++]\n#include <iostream>\n#include <cmath>\nusing namespace std;\n\n// Funcion para calcular area de un circulo\ndouble calcularAreaCirculo(double radio) {\n    return 3.14159 * radio * radio;\n}\n\n// Funcion para calcular area de un rectangulo\ndouble calcularAreaRectangulo(double largo, double ancho) {\n    return largo * ancho;\n}\n\n// Funcion para calcular area de un triangulo\ndouble calcularAreaTriangulo(double base, double altura) {\n    return (base * altura) / 2.0;\n}\n\n// Funcion para mostrar menu de opciones\nvoid mostrarMenu() {\n    cout << \"\\nCalculadora de Areas\" << endl;\n    cout << \"1. Circulo\" << endl;\n    cout << \"2. Rectangulo\" << endl;\n    cout << \"3. Triangulo\" << endl;\n    cout << \"0. Salir\" << endl;\n    cout << \"Elija una opcion: \";\n}\n\nint main() {\n    int opcion;\n    double resultado;\n    \n    do {\n        mostrarMenu();\n        cin >> opcion;\n        \n        if (opcion == 1) {\n            double radio;\n            cout << \"Ingrese el radio del circulo: \";\n            cin >> radio;\n            resultado = calcularAreaCirculo(radio);\n            cout << \"Area del circulo: \" << resultado << endl;\n        }\n        else if (opcion == 2) {\n            double largo, ancho;\n            cout << \"Ingrese largo y ancho del rectangulo: \";\n            cin >> largo >> ancho;\n            resultado = calcularAreaRectangulo(largo, ancho);\n            cout << \"Area del rectangulo: \" << resultado << endl;\n        }\n        else if (opcion == 3) {\n            double base, altura;\n            cout << \"Ingrese base y altura del triangulo: \";\n            cin >> base >> altura;\n            resultado = calcularAreaTriangulo(base, altura);\n            cout << \"Area del triangulo: \" << resultado << endl;\n        }\n        else if (opcion == 0) {\n            cout << \"Saliendo del programa...\" << endl;\n        }\n        else {\n            cout << \"Opcion invalida.\" << endl;\n        }\n    } while (opcion != 0);\n    \n    return 0;\n}\n\\end{lstlisting}\n\nEste ejemplo demuestra cómo las funciones permiten:\n\n\\begin{itemize}\n    \\item \\textbf{Reutilización:} Las funciones de cálculo pueden usarse en múltiples partes del programa\n    \\item \\textbf{Organización:} El código está dividido en tareas específicas\n    \\item \\textbf{Legibilidad:} El flujo del programa es fácil de seguir\n    \\item \\textbf{Mantenibilidad:} Si necesitas cambiar un cálculo, solo lo haces en un lugar\n    \\item \\textbf{Pruebas:} Puedes probar cada función por separado\n\\end{itemize}\n\n\\textbf{Recuerda:} Las funciones son las herramientas fundamentales que te permiten construir programas grandes y complejos dividiéndolos en partes manejables. Son esenciales para escribir código limpio, reutilizable y mantenible.\n